\documentclass{article}

\usepackage[margin=1in]{geometry}
\usepackage{amsmath,amssymb}
\usepackage{algorithm}
\usepackage{algpseudocode}
\usepackage{booktabs}
\usepackage{enumitem}
\usepackage{hyperref}
\usepackage{xcolor}

\newcommand{\Hyp}{\mathcal{H}}
\newcommand{\Conf}{\mathrm{conf}}
\newcommand{\Type}{\mathrm{type}}
\newcommand{\Dep}{\mathrm{dep}}

\title{Memory-Guided Hypothesis Search:\\
A Neuro-Symbolic Framework for Test-Time Reasoning with\\
Multimodal LLMs on Interactive Grid Puzzles}
\author{}
\date{}

\begin{document}
\maketitle

\begin{abstract}
We describe a test-time search framework that lets a multimodal LLM (MLLM)
learn the rules of an unfamiliar, partially-observed environment (an
ARC-AGI-3 game) purely from interaction, without gradient updates. The
framework maintains an explicit, typed \emph{hypothesis memory} —
short natural-language claims about action effects, win conditions,
component roles, and causal rules, each carrying a confidence score,
an evidence trail, and dependency edges onto other hypotheses. Rather
than trusting the MLLM's free-form judgments about what just happened,
the framework treats them as \emph{noisy evidence} to be cross-checked
against cheap, deterministic signals computed directly from the pixel
grid (movement deltas, connected components, edge-strip meters, BFS
shortest paths). This turns a single black-box reasoning loop into a
neuro-symbolic pipeline in which the MLLM is invoked only for the
sub-problems that are genuinely underspecified — perception and
semantic labeling — while every sub-problem that is exactly
specifiable from the grid is solved algorithmically. We present the
hypothesis representation, the per-step search loop, the calibration
procedure that grounds action semantics before any free-form reasoning
begins, the dependency-DAG mechanism that orders what gets tested next,
and the hierarchical decomposition used once enough structure is known
to plan directed action sequences.
\end{abstract}

\section{Problem Setting}
\label{sec:setting}

At each timestep $t$ the agent observes one or more grids
$g_t \in \{0,\dots,15\}^{H\times W}$ (a short animation of the most
recent frames) and a set of available actions $A_t$. Action semantics
are \emph{opaque}: the agent knows only symbolic names
(\texttt{ACTION1}, \dots, \texttt{ACTION6}, \texttt{PLACE}), not their
effects, and must discover both \emph{what the actions do} and
\emph{what the win condition is} through interaction alone. The agent
has a fixed action budget and receives only a terminal
\texttt{WIN}/\texttt{GAME\_OVER} signal from the environment — there is
no reward shaping and no privileged access to game logic.

This is a natural setting for test-time hypothesis search: the agent
must build, in-context, a small causal model of the environment and use
it to act, revising the model as evidence accumulates. The central
design question is \emph{what to trust}: MLLM free-form judgments about
noisy pixel observations are unreliable (Section~\ref{sec:verification}),
so the framework is built around minimizing what is asked of the model
and maximizing what is computed.

\section{Hypothesis Representation}
\label{sec:representation}

The agent's belief state is a set of hypotheses
$\Hyp = \{h_1, \dots, h_n\}$, $n \le N_{\max}=16$. Each hypothesis is a
tuple
\[
h = \big(\,\mathrm{text},\ \tau,\ c \in [0,1],\ k,\ E,\ \Dep(h)\,\big)
\]
where $\mathrm{text}$ is a short natural-language claim, $\tau$ is a
\emph{type} (below), $c$ is a scalar confidence, $k$ counts direct
tests, $E$ is an ordered evidence trail of (text, step, sign) triples,
and $\Dep(h) \subseteq \Hyp$ is a set of prerequisite hypotheses.

\subsection{Type taxonomy}
Every hypothesis is assigned one of four types via a lightweight
lexical classifier applied to its text (no model call required):

\begin{center}
\begin{tabular}{@{}lll@{}}
\toprule
Type & Example & Role \\
\midrule
\textsc{Action\_Effect} & ``\texttt{ACTION1} moves the player UP'' & grounds control semantics \\
\textsc{Win\_Condition} & ``Win by reaching the flashing tile'' & defines the goal test \\
\textsc{Component\_Role} & ``The orange block is the player'' & grounds perception \\
\textsc{Causal\_Rule} & ``Stepping on blue rotates the player'' & environment dynamics \\
\bottomrule
\end{tabular}
\end{center}

The type of a hypothesis determines how it is treated by the rest of
the system: \textsc{Action\_Effect} and player-identity
\textsc{Component\_Role} hypotheses are protected from later
degradation (Section~\ref{sec:lifecycle}); \textsc{Win\_Condition}
hypotheses are held to a stricter confirmation standard (necessary vs.\
sufficient, Section~\ref{sec:verification}); and \textsc{Causal\_Rule}
hypotheses are what the hierarchical planner (Section~\ref{sec:chain})
looks up to decide where to navigate.

\subsection{Confidence dynamics}
Confidence updates asymmetrically:
\[
c \;\leftarrow\;
\begin{cases}
\min(1, c + \delta^+), & \delta^+ = 0.25 & \text{(confirm)}\\
\max(0, c - \delta^-), & \delta^- = 0.15 & \text{(contradict)}
\end{cases}
\]
Confirmation moves confidence faster than contradiction erodes it. This
is deliberate: in this environment a single clean confirmation (e.g.\ a
grid diff that exactly matches a predicted movement direction) is
strong evidence, whereas an apparent contradiction is more often a
confound — the player was blocked by a wall, or a bar/tile effect
muddied the diff — than a genuinely false hypothesis. Hypotheses with
$c < 0.10$ are dropped (Section~\ref{sec:lifecycle}), unless protected.

\subsection{Dependency DAG}
A hypothesis may declare $\Dep(h)$: other hypotheses that must hold
before $h$ can be meaningfully tested (e.g.\ a win-condition hypothesis
that presupposes a specific causal rule). This turns $\Hyp$ into a DAG
over which two queries are defined:

\begin{itemize}[nosep]
\item \textbf{Frontier}$(\Hyp) = \{h : c(h) < 0.80 \ \wedge\ \forall d \in \Dep(h),\ c(d) \ge 0.80\}$ — hypotheses that are not yet confirmed but whose prerequisites are, i.e.\ what is ready to test \emph{next}.
\item \textbf{BlockingAncestor}$(h) = \arg\min_{a \in \mathrm{Ancestors}(h),\, c(a)<0.80} c(a)$ — if $h$ itself is untested but has an unconfirmed prerequisite, this identifies which prerequisite to test first instead.
\end{itemize}

This ordering matters in practice: a downstream goal hypothesis (``win
by reaching the goal'') is untestable until an upstream mechanic
hypothesis (``rotation tiles change facing'') is established, and
testing the goal condition before that is confirmed produces
uninformative, contradictory-looking evidence.

\section{Calibration: Grounding Action Semantics Before Reasoning}
\label{sec:calibration}

Before any free-form LLM reasoning occurs, the agent runs a one-time,
model-light \emph{calibration} pass: each playable action is executed
exactly once, and the resulting pixel diff is compared programmatically
(deterministic grid-diff over the raw grids). The MLLM is asked the same
question (``which block moved, and in which direction?'') only as a
secondary signal, used solely when the deterministic diff is
inconclusive — the programmatic answer takes precedence whenever both
are available.

Two ACTION\_EFFECT / COMPONENT\_ROLE hypotheses are then seeded directly
at confidence $0.95$ and marked \emph{calibration-locked}:

\begin{enumerate}[nosep]
\item \textbf{Action semantics.} Actions are assigned to canonical
UP/DOWN/LEFT/RIGHT slots by fixed convention rather than by trusting the
MLLM's per-action direction label — empirically, MLLM-reported
directions during calibration are unreliable even though grid diffs are
not, so the enumerable, small hypothesis space (4 permutations of 4
directions) is resolved structurally instead of by model judgment.
\item \textbf{Player identity.} The set of colors that shifted position
consistently across \emph{every} calibration action is computed by
set intersection; this is the player sprite. If the diff signal is
degenerate (e.g.\ a click-only game with no directional actions), the
most frequent multi-color MLLM block description is used as a
fallback.
\end{enumerate}

Hypotheses seeded by calibration are exempt from later contradiction or
rewrite by free-form model output (Section~\ref{sec:lifecycle}): they
encode facts established by direct, repeatable observation rather than
by a single ambiguous judgment call, and are treated as ground truth
for the remainder of the run (persisting across levels of the same
game).

\section{The Per-Step Search Loop}
\label{sec:loop}

At each step the agent conceptually makes up to four model calls —
\textsc{Logic}, \textsc{Spatial}, \textsc{Action}, \textsc{Update} — but
in practice most of these are short-circuited by cheaper mechanisms
whenever they apply (Algorithm~\ref{alg:step}).

\begin{algorithm}
\caption{One environment step}
\label{alg:step}
\begin{algorithmic}[1]
\State observe grids $g_t$, available actions $A_t$
\If{previous action exists}
  \State \Call{ProgrammaticVerify}{$g_{t-1}, g_t, \Hyp$} \Comment{\S\ref{sec:verification}}
  \If{\Call{CanSkipLLMUpdate}{$\Hyp$}} skip \textsc{Update} call
  \Else{ } issue \textsc{Update} call, cross-validate against \S\ref{sec:verification} overrides
  \EndIf
  \State \Call{GenerateHypothesesFromObservation}{$g_t$} \Comment{programmatic pattern detectors}
  \State \Call{Prune}{$\Hyp$}
\EndIf
\State $a \gets$ \Call{StrategicActionCheck}{$g_t, \Hyp$} \Comment{\S\ref{sec:chain}, tried first}
\If{$a = \bot$}
  \State $\ell \gets$ \textsc{Logic}$(\Hyp, g_t)$ \Comment{free-form fallback reasoning}
  \State $a \gets$ \textsc{Action}$(\Hyp, g_t, \ell,\ \text{test-directive from \S\ref{sec:representation}.3})$
\EndIf
\State execute $a$; record trajectory
\end{algorithmic}
\end{algorithm}

\paragraph{\textsc{Logic}.} Free-form reasoning over the formatted
hypothesis memory and current frame. This is the least constrained
call and, empirically, the least reliable across MLLM backends tested
(both a 7B image-native model and a mid-size text model showed
frequent malformed or empty output on this call specifically); it is
retained as a fallback narrative for the \textsc{Action} prompt rather
than as a load-bearing reasoning step.

\paragraph{\textsc{Spatial} (optional).} Where supported, predicts the
next grid state. Two modes are available, tried in order of
determinism: (1) parse an explicit predicted grid from the model's text
and render it directly — exact and game-state-aware; (2) for backends
that support image-conditioned generation, condition on the current
frame. Unconditioned text-to-image generation is deliberately never
used as a fallback, since an unconditioned sample is pure noise with
respect to the actual next state and would be worse than omitting the
prediction.

\paragraph{\textsc{Action}.} Selects the next action. The prompt is
composed from: the hypothesis memory, an exploration bonus that
prioritizes untried actions, a DAG-guided test directive
(Section~\ref{sec:representation}.3) naming which under-tested
hypothesis to probe, and — if any \textsc{Causal\_Rule} hypothesis is
confirmed — a strategic directive naming a concrete navigation goal
instead of a generic ``explore'' instruction. Parsing failures degrade
through three levels: full prompt $\to$ a bare enumerated-choice
re-prompt $\to$ a deterministic exploratory pick (least-tried untried
action).

\paragraph{\textsc{Update}.} Judges whether the previous action
confirmed or contradicted each relevant hypothesis. This call is the
one most exposed to hallucination and is treated accordingly — see
Section~\ref{sec:verification}.

\section{Neuro-Symbolic Verification}
\label{sec:verification}

The central reliability mechanism is that \textbf{the MLLM's verdicts
are evidence, not ground truth}. Every hypothesis update is filtered
through deterministic checks computed from the raw grids before being
applied:

\begin{itemize}[nosep]
\item \textbf{Grid-diff movement detection.} A directional
\textsc{Action\_Effect} hypothesis is confirmed/contradicted by
comparing the detected player movement delta between frames against the
direction named in the hypothesis text — with no model call. If the
grid did not change, the system first checks whether the player was
simply \emph{blocked} (wall/boundary) before contradicting, since a
blocked-but-correct direction must not be penalized as a wrong
hypothesis.
\item \textbf{Position-tracking fallback.} When a raw grid diff is
ambiguous (e.g.\ a bar/meter or remote tile effect changed unrelated
cells in the same step), the system falls back to tracked player
position rather than the diff, which is a more robust signal when
multiple things changed at once.
\item \textbf{Edge-strip meter tracking.} Bar/timer/health hypotheses
are verified by comparing colored-cell counts in the frame's edge
strips before and after the action, independent of any model judgment.
\item \textbf{Remote-change detection.} Non-player cells changing
elsewhere on the grid (e.g.\ a tile that rotates a distant pattern) are
detected by diffing outside the player's neighborhood, and used to
auto-generate new \textsc{Causal\_Rule} hypotheses (Section~\ref{sec:lifecycle}) even when the LLM never mentions the effect.
\item \textbf{Skip-LLM shortcut.} If every open hypothesis is either an
\textsc{Action\_Effect}/\textsc{Component\_Role} type or already at
$c \ge 0.90$, the entire \textsc{Update} call is skipped and only the
programmatic checks above run — there is nothing left for free-form
judgment to usefully adjudicate.
\end{itemize}

On top of these, several \emph{guard rules} correct or reject model
verdicts that are structurally impossible given what has already been
established:
\begin{enumerate}[nosep]
\item Persistent identity facts (``X is the player'') are never
weakened or dropped, regardless of verdict.
\item Calibration-locked \textsc{Action\_Effect} hypotheses
(Section~\ref{sec:calibration}) cannot be revised by later free-form
output.
\item A model \textsc{Confirm} of a \textsc{Win\_Condition} hypothesis
is overridden to \textsc{Contradict} if the level did not actually
advance — the condition tested is necessary but evidently not
sufficient — \emph{unless} the hypothesis describes a multi-step
conditional sequence not yet completed.
\item A model \textsc{Contradict} of a directional
\textsc{Action\_Effect} hypothesis is cross-checked against tracked
player position; if the position in fact shifted in the predicted
direction, the verdict is flipped back to \textsc{Confirm} and the
model's stated reason is replaced with the programmatic one.
\end{enumerate}

This design follows directly from an observed failure mode: when the
free-form \textsc{Logic}/\textsc{Update} calls fail or hallucinate
(observed with both an image-native 7B model and a larger text-only
model), a naive pipeline would fall through to an ungrounded, generic
action prompt and visibly oscillate (e.g.\ alternating UP/DOWN). Routing
verification and, where possible, action selection itself
(Section~\ref{sec:chain}) through deterministic grid computation removes
this failure mode structurally rather than by prompting harder.

\section{Hierarchical Strategic Execution}
\label{sec:chain}

Once at least one \textsc{Causal\_Rule} hypothesis is confirmed, a
single open-ended ``decide what to do'' call is replaced by a
\emph{chain of narrow, independently-checkable sub-decisions}
(Algorithm~\ref{alg:chain}). Each link is either a single-purpose model
query with a small, checkable answer space, or a deterministic
computation — never a joint, unconstrained decision.

\begin{algorithm}
\caption{Strategic action check (chain decomposition)}
\label{alg:chain}
\begin{algorithmic}[1]
\Require confirmed rule $r \in \Hyp$, current grid $g_t$
\State $\mathrm{met} \gets$ \textsc{ConditionCheck}$(g_t)$ \Comment{VLM: binary yes/no query}
\If{$\mathrm{met}$}
  \State navigate toward the terminal/goal region; \Return
\EndIf
\State $\mathrm{tile} \gets$ \Call{ExtractTileName}{$r.\text{text}$} \Comment{programmatic string extraction}
\If{stored plan targets $\mathrm{tile}$}
  \State pop next step from stored plan
\Else
  \State $\mathrm{plan} \gets$ \Call{BFSShortestPath}{$g_t$, player pos, $\mathrm{tile}$} \Comment{deterministic}
  \If{BFS fails to resolve target/path}
    \State \Return $\bot$ \Comment{fall through to full LLM pipeline}
  \EndIf
\EndIf
\State $\mathrm{action} \gets$ \Call{DirectionToAction}{direction} \Comment{from calibrated semantics, \S\ref{sec:calibration}}
\State \Return $\mathrm{action}$
\end{algorithmic}
\end{algorithm}

The key design point is the \emph{ordering of trust}: a deterministic
BFS over the parsed grid is attempted first, and a VLM-based direction
estimate is used only as a last resort when BFS cannot confidently
resolve the target tile or a path to it. In earlier iterations, the
order was reversed (VLM direction estimate primary), and it produced
degenerate behavior — voting for a wall-blocked direction repeatedly —
because the visual arithmetic implicit in ``which way is the goal
relative to me'' is exactly the kind of quantity a shortest-path search
computes exactly and a VLM estimates noisily. The general principle:
\textbf{whenever a sub-decision is exactly specifiable from the grid,
compute it; reserve the model call for the sub-decision that is
genuinely a perception/semantics question} (is the condition visually
met, which named tile does the rule refer to).

\section{Hypothesis Lifecycle Management}
\label{sec:lifecycle}

\begin{description}[nosep]
\item[Add] New hypotheses are parsed from \textsc{Logic}/\textsc{Update}
text or generated programmatically from detected patterns (bar shrink,
remote change, panel/legend structure). Near-duplicate text (including
short-form action–direction matches, e.g.\ ``\texttt{ACTION1} UP'' vs.\
an existing ``\texttt{ACTION1} moves the player one cell UP'') is
deduplicated before insertion.
\item[Confirm / Contradict] Apply the asymmetric update of
Section~\ref{sec:representation}.2 and append an evidence record.
\item[Rewrite] Replaces the text of a hypothesis whose framing was
wrong but whose target phenomenon was real; resets confidence to a
neutral $0.4$ and clears prior evidence, since old evidence no longer
applies to the new claim.
\item[Drop] Removes a hypothesis once $c < 0.10$, except for
protected types (player identity, calibration-locked
\textsc{Action\_Effect}). Dependents lose the dropped id from their
$\Dep$ set rather than being recursively invalidated.
\item[Evict] When $|\Hyp| = N_{\max}$, the lowest-confidence hypothesis
is evicted to make room for a new one, drawing preferentially from
non-\textsc{Action\_Effect} hypotheses — calibrated action knowledge is
never sacrificed for capacity.
\end{description}

\section{Design Principles}
\label{sec:principles}

The mechanisms above follow from a small number of repeated design
choices, stated here for reuse in other test-time-search agents built
on unreliable or small MLLMs:

\begin{enumerate}[nosep]
\item \textbf{Compute what is exactly specifiable; ask the model only
what is genuinely underspecified.} Movement, shortest paths, connected
components, and meter levels are computed from the grid; only
perception (``what does this shape represent'') and semantics
(``is the visually-apparent condition satisfied'') go to the model.
\item \textbf{Treat model verdicts as evidence, not ground truth.}
Every confirm/contradict is filtered through cheap, deterministic
cross-checks and a small set of structural guard rules before being
applied to the belief state.
\item \textbf{Decompose compound decisions into narrow, checkable
sub-queries.} A single open-ended ``what should I do'' call is
replaced by a chain of single-purpose calls (condition check, rule
lookup, locate, act), each with a small answer space that is easy to
validate or override.
\item \textbf{Establish irreversible facts early and lock them.}
Calibration determines action semantics and player identity once,
deterministically, before any free-form reasoning occurs, and these
facts are then protected from later corruption by noisy model output.
\item \textbf{Order hypothesis testing by dependency, not by
recency or confidence alone.} A DAG over hypotheses ensures
prerequisite mechanics are established before downstream goal
conditions are tested against them, avoiding uninformative or
misleading contradictions.
\end{enumerate}

\end{document}
